\documentclass[12pt,oneside,noprintercorrection]{iut}

% Affichage des aides à la rédaction
\newif\ifshow
% Commenter la ligne suivante pour faire disparaitre les éléments d'aide à la rédaction
\showtrue

\usepackage{macros}

%-------------------------------------------------------------------
%                         Page de titre:
%-------------------------------------------------------------------



\ThesisTitle{Titre du rapport}
\ThesisKind{Rapport de mission}
\ThesisPresentedThe{soutenu le XXXX}
\ThesisAuthor{Prénom NOM}

\NomDuLaboOuEntreprise{Nom de l'entreprise}
\LogoLaboOuEntreprise{img/blankImage} % Image du logo du labo ou ets dans le rep img

\EncadrantEts = {XXX}
\EncadrantUniv = {YYY}

\begin{document}

% Creation de la page de titre:
\MakeThesisTitlePage



%-------------------------------------------------------------------


%-------------------------------------------------------------------
%                          remerciements
%-------------------------------------------------------------------

%\DontFrameThisInToc
\begin{ThesisAcknowledgments}

\end{ThesisAcknowledgments}





%-------------------------------------------------------------------
%                  ecriture de `Chapitre' et `Partie'
%                      dans la table des matieres
%-------------------------------------------------------------------

\WritePartLabelInToc \WriteChapterLabelInToc


%-------------------------------------------------------------------
%                        table des matieres
%-------------------------------------------------------------------

\tableofcontents

%-------------------------------------------------------------------
%              Exemple d'utilisation de \SpecialSection
%-------------------------------------------------------------------

% La commande \mainmatter (nouvelle commande LaTeX2e) permet de passer
% a la numerotation arabe (ce que fait \pagenumbering{arabic})
% et de faire commencer la nouvelle page 1 sur une page impaire.
% On evitera donc d'utiliser directement \pagenumbering{arabic}.
\mainmatter

% ----------------------------------------------------------------
\SpecialSection{Résumés}

Résumés en français et en anglais du rapport.

\SpecialSection{Introduction}

\begin{aide}{}{}
  \begin{itemize}
    \item Commentez \verb!\showtrue! ligne 6 pour faire disparaitre les éléments d'aide à la rédaction.
    \item Le nombre de pages doit être compris entre 15 et 30 pages, sans les annexes.
    \item Il n'y a pas d'obligation d'utiliser \LaTeX pour écrire votre rapport. Vous pouvez utiliser un autre logiciel mais la mise en page finale doit être la plus proche possible de ce modèle.
  \end{itemize}
\end{aide}

\begin{aide}{}{}
  L'introduction doit comprendre 2 à 4 pages. En plus d'informations factuelles comme l'année d'étude (BUT 2 ou 3), on doit y trouver une description de l'entreprise ainsi qu'une description du contexte de la mission. Quelle est l'importance et l'utilité de la mission pour l'entreprise ? Quel était l'existant par rapport à la mission ? (nouveau projet, aucun code existant ; code déjà existant \dots)

  L'introduction doit ensuite donner une vue d'ensemble de la mission du stage et terminer par une annonce du plan.

\end{aide}

\begin{aide}{}{}
  Le coeur du rapport doit typiquement comporter trois chapitres qui vont mettre en valeur différents aspects de votre année et démontrer votre prise de recul par rapport à ce que vous avez fait. Le rapport n'a pas nécessairement besoin de suivre chronologiquement ce qui a été fait durant cette année. Il faut avant tout faire une présentation logique de votre travail. Réfléchissez aux différentes missions réalisées. Quelle est la plus importante ? Si vous avez travaillé sur une seule mission, quelle a été la partie la plus importante ? Un point technique doit être abordé dans votre rapport : qu'est-ce qui a été le plus difficile au niveau technique ?

  Il est conseillé :
  \begin{itemize}
    \item  d'utiliser des diagrammes (un MCD pour une Base de données, un Diagramme de classes pour une appli java \dots) pour illustrer le fonctionnement global d'une application ou d'un système. Ces diagrammes sont souvent utiles pour la compréhension~;
    \item de ne pas mettre de logos de logiciels (exemple celui d'Eclipse). Cela n'apporte rien;
    \item d'utiliser des copies d'écran en nombre restreint et pour illustrer des choses difficiles à expliquer simplement avec du texte. Par exemple, quand une interface a été développée, il peut être utile d'afficher une capture d'écran pour la montrer. N'utilisez pas de copies d'écran pour montrer du code~;
    \item d'apporter un soin particulier à l'orthographe et la grammaire. Vérifiez que toutes vos phrases ont du sens. Relisez-vous et faites relire.
  \end{itemize}

\end{aide}

\begin{aide}{}{}
Toute forme de plagiat est proscrite et sera sévèrement sanctionnée si elle est découverte.
\end{aide}

\begin{aide}{}{}
  Pour les BUT 3, il est attendu de réaliser la présentation d'un outil ou d'une technologie que vous avez appris ou qui vous semble important durant votre année, sous la forme d'un mini cours de 2 à 3 pages.
  L'objectif de ce mini-cours est de permettre à un lecteur ne connaissant pas cet outil ou technique de pouvoir le mettre en oeuvre sur un exemple simple. On pourra préciser notamment :
  \begin{itemize}
  \item Qu'est ce que c'est ?
  \item A quoi ça sert ?
  \item Comment on le met en oeuvre (qu'est-ce qu'il faut installer ?).
  \item Un exemple trivial de mise en oeuvre (comment bien démarrer).
  \end{itemize}
  Vous privilégierez bien évidemment, quand c'est possible, des outils ou techniques qui n'ont pas été étudiées durant vos études.
\end{aide}

% Pour ne pas avoir le mot `Chapitre' au debut de chaque chapitre.
\NoChapterHead


%--------------------------------------------------------------------------------------
%--------------------------------------------------------------------------------------
\chapter{Prise en main de \LaTeX}

Il est possible d'utiliser {\LaTeX} directement dans un navigateur avec des outils comme OverLeaf mais vous pouvez aussi compiler le code {\LaTeX} directement sur votre machine.

\section{Installation}

\subsection{Sous Linux}
Nombreuses documentations disponibles sur internet pour
l'installation des packages. Sous {\sc Ubuntu} par exemple, 
le package {\tt texlive} installe une sélection des outils les 
plus fréquements utilisés.

\subsection{Sous Windows}

\begin{itemize}
\item Pour compiler les fichiers {\tt .tex} en {\tt .pdf}, installer  {\sc Miktex} \cite{Miktex}
\item Pour écrire des documents \LaTeX, installer  {\sc TeXnicenter} \cite{Texnicenter} ou  {\sc ConTEXT} \cite{ConTEXT}.
\end{itemize}

\subsection{Sous Mac}
Utiliser par exemple MacTeX. 

\subsection{Figures}
\label{subsection:figures}

Pour créer vos propres figures, vous pouvez utiliser {\sc Inkscape} \cite{Inkscape} (ou éventuellement l'outil {\sc Draw} d'OpenOffice \cite{OpenOffice}) pour réaliser des dessins vectoriels; il est également possible d'utiliser {\sc Gimp} \cite{Gimp} pour réaliser des dessins bitmaps.

Notez qu'{\sc Inkscape} sait réaliser la plupart des conversions vectorielles vers {\tt pdf}, et que {\sc Gimp} sait réaliser la plupart des conversions bitmaps vers {\tt png} ($\rightarrow$ schémas) ou {\tt jpg} ($\rightarrow$ photos).

\section{Compilation des documents}
Pour compiler un document \LaTeX  en pdf, le plus simple est d'utiliser la commande {\tt pdflatex}; 
il est pour cela nécessaire d'inclure les images dans les formats {\tt .pdf}, {\tt .jpg}, ou {\tt .png}
comme indiqué en section \ref{subsection:figures}.


Par exemple, pour compiler ce document, les commandes suivantes ont été lancées :
\begin{center}
  \fbox{
    {\scriptsize\tt
      pdflatex rapport~;~~ bibtex   rapport~;~~pdflatex rapport~;~~pdflatex rapport
    }
  }
\end{center}            
\section{Quelques commandes}
\subsection{Insertion de figures}

\begin{figure}[!ht]
  \centering
  \includegraphics[width=8cm]{img/UL}
  \caption{Le logo de l'Université de Lille.}
  \label{fig:univ-lille}
\end{figure}

Voici le Logo de l'Université de Lille (voir fig. \ref{fig:univ-lille}) au format vectoriel. Vous pouvez mettre les images dans le répertoire \texttt{img} si vous modifiez ce fichier exemple ou en ajoutez d'autres selon ce modèle.

\subsection{Insertion d'équations}

Les équations (et autres formules) sont un des points forts de \LaTeX, utile si vous devez formaliser votre travail. Quelques exemples ci-dessous\ldots

Pour une fonction $f(x)$ continue et croissante sur l'intervalle $[a..b]$, l'équation \ref{eq:sumint} sert à \ldots

\begin{equation}
\sum_{i=a}^{b-1} f(i)   \leq   \int_{a}^{b}f(t)\;dt   \leq   \sum_{i=a+1}^{b} f(i)
\label{eq:sumint}
\end{equation}

ou alors la matrice Vandermonde \ref{eq:matvan} sert à \ldots

\begin{equation}
V=
\left( 
\begin{array}{ccccc}
1 & \alpha_1 & \alpha_1^2 & \dots & \alpha_1^{n-1}\\
1 & \alpha_2 & \alpha_2^2 & \dots & \alpha_2^{n-1}\\
1 & \alpha_3 & \alpha_3^2 & \dots & \alpha_3^{n-1}\\
\vdots & \vdots & \vdots & \ddots &\vdots \\
1 & \alpha_m & \alpha_m^2 & \dots & \alpha_m^{n-1}\\
\end{array}
\right) 
\label{eq:matvan}
\end{equation}

ou encore un définition récursive \ref{eq:fac} qui peut servir dans certains cas à \ldots

\begin{equation} 
  fact(n) = \left\{ 
    \begin{array}{ll}
      1                  & \mbox{si $x \leq 1$};\\
      n \times fact(n-1) & \mbox{autrement}.
    \end{array} 
  \right.
  \label{eq:fac}
\end{equation}


Un document qui peut vous être utile est le suivant \cite{mshort}. La documentation {\em The Not So Short Introduction to \LaTeX} \cite{lshort} présente également des exemples mathématique assez détaillés.
\newpage % saut de page pour éviter certains mauvais formatages

\subsection{Insertion de code}

Pour plus d'infos sur le package {\tt listings}, consulter cette note de bas de page\footnote{
\url{ftp://tug.ctan.org/pub/tex-archive/macros/latex/contrib/listings/listings.pdf}

}

\begin{lstlisting}[frame=trBL]
#include <iostream>

int main() {
    std::cout << "Hello, world!\n";
}
\end{lstlisting}

\section{Doc \LaTeX}
Faire des recherches sur Google ou consulter ce livre très complet \cite{LatexCompanion}. La documentation {\em The Not So Short Introduction to \LaTeX} \cite{lshort} est également un très bon point de départ; elle est disponible en ligne.

\SpecialSection{Conclusion}

\begin{aide}{}{}
  La conclusion doit dresser un bilan de votre année et récapituler les contributions (apports des missions pour l'entreprise, le client\dots). C'est aussi le lieu pour discuter des difficultés rencontrées et des moyens utilisés pour les lever. Enfin, vous dresserez un bilan humain et personnel en termes de savoir faire et savoir être. Cette année en alternance vous a-t-elle aidé à choisir ou confirmer votre projet personnel ?

  La conclusion fait typiquement 1 à 2 pages.
\end{aide}

\begin{aide}{}{}
Pour les BUT 3, il est demandé de réaliser une courte comparaison entre le stage/l'alternance de BUT 2 et l'année en entreprise de BUT 3.
\end{aide}

\bibliographystyle{myunsrt}
\small
\bibliography{rapport}

\Annex{Annexe 1}

\begin{aide}{}{}
  Les annexes ne sont généralement pas nécessaires sauf cas particulier.
\end{aide}

\end{document}
